% Options for packages loaded elsewhere
% Options for packages loaded elsewhere
\PassOptionsToPackage{unicode}{hyperref}
\PassOptionsToPackage{hyphens}{url}
\PassOptionsToPackage{dvipsnames,svgnames,x11names}{xcolor}
%
\documentclass[
  letterpaper,
  DIV=11,
  numbers=noendperiod]{scrreprt}
\usepackage{xcolor}
\usepackage{amsmath,amssymb}
\setcounter{secnumdepth}{5}
\usepackage{iftex}
\ifPDFTeX
  \usepackage[T1]{fontenc}
  \usepackage[utf8]{inputenc}
  \usepackage{textcomp} % provide euro and other symbols
\else % if luatex or xetex
  \usepackage{unicode-math} % this also loads fontspec
  \defaultfontfeatures{Scale=MatchLowercase}
  \defaultfontfeatures[\rmfamily]{Ligatures=TeX,Scale=1}
\fi
\usepackage{lmodern}
\ifPDFTeX\else
  % xetex/luatex font selection
\fi
% Use upquote if available, for straight quotes in verbatim environments
\IfFileExists{upquote.sty}{\usepackage{upquote}}{}
\IfFileExists{microtype.sty}{% use microtype if available
  \usepackage[]{microtype}
  \UseMicrotypeSet[protrusion]{basicmath} % disable protrusion for tt fonts
}{}
\makeatletter
\@ifundefined{KOMAClassName}{% if non-KOMA class
  \IfFileExists{parskip.sty}{%
    \usepackage{parskip}
  }{% else
    \setlength{\parindent}{0pt}
    \setlength{\parskip}{6pt plus 2pt minus 1pt}}
}{% if KOMA class
  \KOMAoptions{parskip=half}}
\makeatother
% Make \paragraph and \subparagraph free-standing
\makeatletter
\ifx\paragraph\undefined\else
  \let\oldparagraph\paragraph
  \renewcommand{\paragraph}{
    \@ifstar
      \xxxParagraphStar
      \xxxParagraphNoStar
  }
  \newcommand{\xxxParagraphStar}[1]{\oldparagraph*{#1}\mbox{}}
  \newcommand{\xxxParagraphNoStar}[1]{\oldparagraph{#1}\mbox{}}
\fi
\ifx\subparagraph\undefined\else
  \let\oldsubparagraph\subparagraph
  \renewcommand{\subparagraph}{
    \@ifstar
      \xxxSubParagraphStar
      \xxxSubParagraphNoStar
  }
  \newcommand{\xxxSubParagraphStar}[1]{\oldsubparagraph*{#1}\mbox{}}
  \newcommand{\xxxSubParagraphNoStar}[1]{\oldsubparagraph{#1}\mbox{}}
\fi
\makeatother


\usepackage{longtable,booktabs,array}
\newcounter{none} % for unnumbered tables
\usepackage{calc} % for calculating minipage widths
% Correct order of tables after \paragraph or \subparagraph
\usepackage{etoolbox}
\makeatletter
\patchcmd\longtable{\par}{\if@noskipsec\mbox{}\fi\par}{}{}
\makeatother
% Allow footnotes in longtable head/foot
\IfFileExists{footnotehyper.sty}{\usepackage{footnotehyper}}{\usepackage{footnote}}
\makesavenoteenv{longtable}
\usepackage{graphicx}
\makeatletter
\newsavebox\pandoc@box
\newcommand*\pandocbounded[1]{% scales image to fit in text height/width
  \sbox\pandoc@box{#1}%
  \Gscale@div\@tempa{\textheight}{\dimexpr\ht\pandoc@box+\dp\pandoc@box\relax}%
  \Gscale@div\@tempb{\linewidth}{\wd\pandoc@box}%
  \ifdim\@tempb\p@<\@tempa\p@\let\@tempa\@tempb\fi% select the smaller of both
  \ifdim\@tempa\p@<\p@\scalebox{\@tempa}{\usebox\pandoc@box}%
  \else\usebox{\pandoc@box}%
  \fi%
}
% Set default figure placement to htbp
\def\fps@figure{htbp}
\makeatother

\ifLuaTeX
  \usepackage{luacolor}
  \usepackage[soul]{lua-ul}
\else
  \usepackage{soul}
\fi

% definitions for citeproc citations
\NewDocumentCommand\citeproctext{}{}
\NewDocumentCommand\citeproc{mm}{%
  \begingroup\def\citeproctext{#2}\cite{#1}\endgroup}
\makeatletter
 % allow citations to break across lines
 \let\@cite@ofmt\@firstofone
 % avoid brackets around text for \cite:
 \def\@biblabel#1{}
 \def\@cite#1#2{{#1\if@tempswa , #2\fi}}
\makeatother
\newlength{\cslhangindent}
\setlength{\cslhangindent}{1.5em}
\newlength{\csllabelwidth}
\setlength{\csllabelwidth}{3em}
\newenvironment{CSLReferences}[2] % #1 hanging-indent, #2 entry-spacing
 {\begin{list}{}{%
  \setlength{\itemindent}{0pt}
  \setlength{\leftmargin}{0pt}
  \setlength{\parsep}{0pt}
  % turn on hanging indent if param 1 is 1
  \ifodd #1
   \setlength{\leftmargin}{\cslhangindent}
   \setlength{\itemindent}{-1\cslhangindent}
  \fi
  % set entry spacing
  \setlength{\itemsep}{#2\baselineskip}}}
 {\end{list}}
\usepackage{calc}
\newcommand{\CSLBlock}[1]{\hfill\break\parbox[t]{\linewidth}{\strut\ignorespaces#1\strut}}
\newcommand{\CSLLeftMargin}[1]{\parbox[t]{\csllabelwidth}{\strut#1\strut}}
\newcommand{\CSLRightInline}[1]{\parbox[t]{\linewidth - \csllabelwidth}{\strut#1\strut}}
\newcommand{\CSLIndent}[1]{\hspace{\cslhangindent}#1}



\setlength{\emergencystretch}{3em} % prevent overfull lines

\providecommand{\tightlist}{%
  \setlength{\itemsep}{0pt}\setlength{\parskip}{0pt}}



 


\KOMAoption{captions}{tableheading}
\makeatletter
\@ifpackageloaded{tcolorbox}{}{\usepackage[skins,breakable]{tcolorbox}}
\@ifpackageloaded{fontawesome5}{}{\usepackage{fontawesome5}}
\definecolor{quarto-callout-color}{HTML}{909090}
\definecolor{quarto-callout-note-color}{HTML}{0758E5}
\definecolor{quarto-callout-important-color}{HTML}{CC1914}
\definecolor{quarto-callout-warning-color}{HTML}{EB9113}
\definecolor{quarto-callout-tip-color}{HTML}{00A047}
\definecolor{quarto-callout-caution-color}{HTML}{FC5300}
\definecolor{quarto-callout-color-frame}{HTML}{acacac}
\definecolor{quarto-callout-note-color-frame}{HTML}{4582ec}
\definecolor{quarto-callout-important-color-frame}{HTML}{d9534f}
\definecolor{quarto-callout-warning-color-frame}{HTML}{f0ad4e}
\definecolor{quarto-callout-tip-color-frame}{HTML}{02b875}
\definecolor{quarto-callout-caution-color-frame}{HTML}{fd7e14}
\makeatother
\makeatletter
\@ifpackageloaded{bookmark}{}{\usepackage{bookmark}}
\makeatother
\makeatletter
\@ifpackageloaded{caption}{}{\usepackage{caption}}
\AtBeginDocument{%
\ifdefined\contentsname
  \renewcommand*\contentsname{Table of contents}
\else
  \newcommand\contentsname{Table of contents}
\fi
\ifdefined\listfigurename
  \renewcommand*\listfigurename{List of Figures}
\else
  \newcommand\listfigurename{List of Figures}
\fi
\ifdefined\listtablename
  \renewcommand*\listtablename{List of Tables}
\else
  \newcommand\listtablename{List of Tables}
\fi
\ifdefined\figurename
  \renewcommand*\figurename{Figure}
\else
  \newcommand\figurename{Figure}
\fi
\ifdefined\tablename
  \renewcommand*\tablename{Table}
\else
  \newcommand\tablename{Table}
\fi
}
\@ifpackageloaded{float}{}{\usepackage{float}}
\floatstyle{ruled}
\@ifundefined{c@chapter}{\newfloat{codelisting}{h}{lop}}{\newfloat{codelisting}{h}{lop}[chapter]}
\floatname{codelisting}{Listing}
\newcommand*\listoflistings{\listof{codelisting}{List of Listings}}
\usepackage{amsthm}
\theoremstyle{definition}
\newtheorem{definition}{Definition}[chapter]
\theoremstyle{remark}
\AtBeginDocument{\renewcommand*{\proofname}{Proof}}
\newtheorem*{remark}{Remark}
\newtheorem*{solution}{Solution}
\newtheorem{refremark}{Remark}[chapter]
\newtheorem{refsolution}{Solution}[chapter]
\makeatother
\makeatletter
\makeatother
\makeatletter
\@ifpackageloaded{caption}{}{\usepackage{caption}}
\@ifpackageloaded{subcaption}{}{\usepackage{subcaption}}
\makeatother
\usepackage{bookmark}
\IfFileExists{xurl.sty}{\usepackage{xurl}}{} % add URL line breaks if available
\urlstyle{same}
\hypersetup{
  pdftitle={Fitted Q-Iteration: Estimación de niveles de presas},
  pdfauthor={Erika Rubí Jiménez Guzmán; Asesor: Dr.~Yofre Hernán García Gómez},
  colorlinks=true,
  linkcolor={blue},
  filecolor={Maroon},
  citecolor={Blue},
  urlcolor={Blue},
  pdfcreator={LaTeX via pandoc}}


\title{Fitted Q-Iteration: Estimación de niveles de presas}
\author{Erika Rubí Jiménez Guzmán \and Asesor: Dr.~Yofre Hernán García
Gómez}
\date{2026-09-18}
\begin{document}
\maketitle

\renewcommand*\contentsname{Table of contents}
{
\setcounter{tocdepth}{2}
\tableofcontents
}

\bookmarksetup{startatroot}

\chapter*{Resumen}\label{resumen}
\addcontentsline{toc}{chapter}{Resumen}

\markboth{Resumen}{Resumen}

Aquí pondré un resumen de mi tesis, que se llama ``Fitted Q-Iteration:
Estimación de niveles de presas''.

\bookmarksetup{startatroot}

\chapter*{Introduction}\label{introduction}
\addcontentsline{toc}{chapter}{Introduction}

\markboth{Introduction}{Introduction}

Esta será la introducción.

ejemplo de una cita Puterman (1994).

\part{Preliminares}

\chapter{Fundamentos Matemáticos}\label{fundamentos-matemuxe1ticos}

\section{Inteligencia Artificial (IA)}\label{inteligencia-artificial-ia}

La \textbf{Inteligencia Artificial (IA)} es un campo de estudio que se
centra en la creación de sistemas capaces de realizar tareas que
normalmente requieren inteligencia humana. La IA combina diversas
disciplinas, incluyendo matemáticas, estadísticas, informática y
neurociencia, para desarrollar algoritmos y modelos que permiten a las
máquinas aprender, razonar y tomar decisiones de manera autónoma.

Este campo está dividido en varias subáreas, cada una con su enfoque y
aplicaciones específicas. Entre las más destacadas se encuentran el
\textbf{Aprendizaje Automático (Machine Learning)}, \textbf{Aprendizaje
Profundo (Deep Learning)}, \textbf{Procesamiento del Lenguaje Natural},
\textbf{Visión por Computadora}, \textbf{Robótica}, entre otras. Cada
una de estas subáreas utiliza diferentes técnicas y enfoques para
abordar problemas complejos y mejorar la capacidad de las máquinas para
interactuar con el mundo que las rodea.

\subsection{Aprendizaje Automático}\label{aprendizaje-automuxe1tico}

En este subcampo de la Inteligencia Artificial, las máquinas aprenden a
partir de datos, similar a como lo hacemos nosotros, permitiendoles
mejorar con la experiencia. El aprendizaje automático se basa en
algoritmos que identifican patrones y relaciones en los datos, lo que
permite a las máquinas hacer predicciones o tomar decisiones sin ser
explícitamente programadas para cada tarea específica. Este enfoque ha
revolucionado diversas industrias, desde la medicina hasta la economía,
al proporcionar herramientas más eficientes y precisas para el análisis
de datos y la toma de decisiones.

El aprendizaje automático se clasifica en tres categorías principales:
\textbf{Aprendizaje Supervisado}, \textbf{Aprendizaje No Supervisado} y
\textbf{Aprendizaje por Refuerzo}, cada una con sus propias técnicas y
aplicaciones. Estas categorías permiten abordar diferentes tipos de
problemas y desafíos, desde la clasificación de datos hasta la
optimización de estrategias en entornos dinámicos.

\begin{definition}[Aprendizaje
Supervisado]\protect\hypertarget{def-sub}{}\label{def-sub}

Esta categoría se basa en utilizar conjuntos de datos etiquetados (datos
de entrada con sus correspondientes salidas) para entrenar modelos
matemáticos que puedan identificar patrones y a partir de ellos predecir
resultados futuros sobre datos nuevos.

Este tipo de aprendizaje es principalmente utilizado en:

\begin{itemize}
\tightlist
\item
  \textbf{\emph{Tareas de clasificación:}} los modelos predicen la clase
  a las que pertenece cada dato. Por ejemplo, identificar si un correo
  electrónico es spam o no.
\item
  \textbf{\emph{Regresión:}} predicen valores numéricos a partir de
  datos de entrada. Por ejemplo, predecir el precio de una vivienda
  basado en sus características.
\end{itemize}

\end{definition}

\begin{definition}[Aprendizaje No
Supervisado]\protect\hypertarget{def-sub}{}\label{def-sub}

A diferencia del aprendizaje supervisado, en este tipo de aprendizaje no
se utilizan datos etiquetados. En su lugar, los algoritmos buscan
patrones ocultos o agrupaciones dentro de los datos. Estos modelos son
utilizados para tareas especificas como:

\begin{itemize}
\tightlist
\item
  \textbf{\emph{Agrupamiento (Clustering):}} los modelos identifican
  grupos de datos similares entre sí. Por ejemplo, segmentar clientes en
  diferentes grupos según sus comportamientos de compra.
\item
  \textbf{\emph{Reducción de dimensionalidad:}} los modelos simplifican
  los datos manteniendo la mayor cantidad de información posible. Por
  ejemplo, reducir el número de variables en un conjunto de datos para
  facilitar su análisis.
\end{itemize}

\end{definition}

\begin{definition}[Aprendizaje por
Refuerzo]\protect\hypertarget{def-sub}{}\label{def-sub}

En este tipo de aprendizaje automático, los agentes autónomos aprenden a
tomar decisiones a partir de la interección con un entorno en el cual
reciben recompensas o penalizaciones según sus acciones. A continuación
se presentan los conceptos fundamentales de este tipo de aprendizaje

\begin{itemize}
\tightlist
\item
  \textbf{\emph{Agente:}} es el sistema que toma decisiones y aprende de
  la experiencia.
\item
  \textbf{\emph{Entorno:}} en donde el agente opera y recibe
  retroalimentación.
\item
  \textbf{\emph{Estado:}} Situación actual del agente en el entorno.
\item
  \textbf{\emph{Acción:}} Decisión que el agente toma en un estado
  determinado.
\item
  \textbf{\emph{Recompensa:}} Retroalimentación que el agente recibe
  después de realizar una acción, indicando si fue beneficiosa o no.
\item
  \textbf{\emph{Política:}} Estrategia que el agente sigue para decidir
  qué acción tomar en cada estado.
\end{itemize}

El aprendizaje por refuerzo se centra esencialmente en la relación entre
un agente, entorno y el objetivo (lo que se desea lograr). La literatura
formula esta relación en términos de un \textbf{Proceso de Decisión de
Markov (MDP)}.

\end{definition}

Véase IBM (2025) para más información sobre estos tipos de aprendizaje y
sus aplicaciones.

\subsubsection{Q-Learning}\label{q-learning}

Existen distintos algoritmos de aprendizaje por refuerzo, entre los
cuales se encuentra el \textbf{Q-Learning}, este permite a un agente
aprender a tomar decisiones en un entorno para maximizar una recompensa
acumulada a lo largo del tiempo. Este modelo es model-free, es decir, no
requiere un modelo del entorno para aprender, sino que se basa en la
experiencia directa del agente a través de la exploración y explotación
de su entorno.

\chapter{Fundamentos Hidráulicos}\label{fundamentos-hidruxe1ulicos}

Algoritmo de aprendizaje por refuerzo\ldots{}

\begin{definition}[Fitted
Q-Iteration]\protect\hypertarget{def-sub}{}\label{def-sub}

Algoritmo Offline (aprende analizando el historial de
experiencias)\ldots{}

\end{definition}

\part{Modelo Matemático}

\chapter{Modelo Matemático para Optimizar la Operación de un Sistema
Hidroeléctrico.}\label{modelo-matemuxe1tico-para-optimizar-la-operaciuxf3n-de-un-sistema-hidroeluxe9ctrico.}

El objetivo de este capítulo es presentar una formulación rigurosa del
problema de optimización de la operación de un sistema hidroeléctrico
como un Proceso de Decisión de Markov (MDP), junto con una descripción
detallada del algoritmo de Fitted Q-Iteration (FQI) utilizado para
aproximar la política óptima. Se enfatizará la conexión entre los
elementos matemáticos formales y su implementación numérica, así como la
validación operativa mediante simulación en lazo cerrado.

\section{Formulación como Proceso de Decisión de Markov
(MDP)}\label{formulaciuxf3n-como-proceso-de-decisiuxf3n-de-markov-mdp}

Siguiendo la notación de Puterman (1994), se define el MDP como la
colección de objetos
\[\{(T, \mathcal{S}, \mathcal{A}, \mathbb{P}, \mathcal{R})\},\] donde
cada componente se especifica a continuación.

\subsection{\texorpdfstring{Estructura Temporal
\((T)\)}{Estructura Temporal (T)}}\label{estructura-temporal-t}

Se considera un horizonte temporal de operación compuesto por \(T = 50\)
años históricos (periodo 1959--2008), en donde cada uno de ellos se
divide en \(q = 24\) quincenas. Con el fin de reducir la dimensionalidad
temporal y capturar la estacionalidad, las quincenas se agrupan en
\(M = 6\) etapas hidrológicas mediante la función de agregación
\(\mu: \{0,\dots,23\} \to \{0,\dots,5\}\) definida por
\begin{equation}\protect\phantomsection\label{eq-etapas}{\mu(q) = \begin{cases} 0, & 0 \leq q <10\\ 1, & 10 \leq q <14\\ 2, & 14 \leq q < 16\\ 3, & 16 \leq q < 18\\ 4, & 18 \leq q < 20\\ 5, & 20 \leq q \leq 23\\ \end{cases}}\end{equation}
de modo que \(m = \mu(q)\) es la etapa hidrológica correspondiente a la
quincena \(q\). Esta partición reproduce exactamente la agrupación
implementada en el código (\texttt{etapas\_quinc}), verificada como
correcta.

\textbf{Dos resoluciones temporales distintas.} Es indispensable
distinguir, desde ahora, dos usos diferentes de esta estructura
temporal: (i) para \textbf{construir el conjunto de aprendizaje} del
algoritmo de FQI (Sec. Section~\ref{sec-conjunto-aprendizaje}), la
afluencia se agrega a nivel de \emph{etapa} (\(M=6\) decisiones por año)
y no se recorre la serie quincena por quincena; (ii) para
\textbf{desplegar y evaluar la política} ya aprendida (Sec.
Section~\ref{sec-simulacion}), la decisión sí se reconsulta en cada una
de las 24 quincenas del año, usando en cada una la afluencia real de esa
quincena, aunque los parámetros vigentes (capacidad máxima, curva guía,
grilla de acción) sean los de la etapa \(m=\mu(q)\) que la contiene.

\subsection{\texorpdfstring{Espacio de estados
\((\mathcal{S})\)}{Espacio de estados (\textbackslash mathcal\{S\})}}\label{sec-espacio-estados}

El sistema está compuesto por dos presas:
\href{https://maps.app.goo.gl/NHNQdHdfYAoRGxhPA}{La Angostura}
\((i = 1)\) y \href{https://maps.app.goo.gl/acX7hVXswX5uYYUBA}{Malpaso}
\((i =2)\). El volumen almacenado en la presa \(i\) en la quincena \(q\)
del año \(t\) se representa por
\(V_{i,t,q} \in \mathbb{R}_{+} \cup \{0\}\) (millones de metros cúbicos,
Mm³).

A diferencia del enfoque de Programación Dinámica Estocástica (PDE), en
el que el almacenamiento se discretiza en celdas y el estado es un
índice entero, \textbf{en FQI el almacenamiento se mantiene como
variable continua}. El parámetro \(\Delta = 600\) Mm³ (el cual permite
subdividir la capacidad de almacenamiento de cada presa) y los enteros
\(N_1 = 27\), \(N_2 = 16\) (capacidades discretas máximas de cada presa)
---heredados de la discretización de la PDE precedente--- no se usan
aquí para discretizar el estado, sino únicamente para fijar las
\textbf{capacidades máximas continuas} de cada presa,
\begin{equation}\protect\phantomsection\label{eq-capacidad-maxima}{V_{1,\max} := N_1 \Delta = 16{,}200 \ \text{Mm}^3, \qquad V_{2,\max} := N_2 \Delta = 9{,}600\ \text{Mm}^3,}\end{equation}
y para convertir a volumen la curva guía, definida originalmente en
unidades de índice (Sec. Section~\ref{sec-recompensa}).

El espacio de estados del MDP es entonces
\begin{equation}\protect\phantomsection\label{eq-espacio-estados}{
\begin{aligned}
\mathcal{S}
&= [0,V_{1,\max}]\times[0,V_{2,\max}]
   \times\{0,\dots,5\} \\
&= \left\{s=(s_1,s_2,m):
s_i\in[0,V_{i,\max}],\ i=1,2,\ 
m\in\{0,\dots,5\}\right\}.
\end{aligned}
}\end{equation}

\textbf{Grilla gruesa de referencia.} Aunque el estado en sí no se
discretiza, tanto la construcción del conjunto de aprendizaje (Sec.
Section~\ref{sec-conjunto-aprendizaje}) como la aproximación del
operador de maximización (Sec. Section~\ref{sec-max-operador}) requieren
evaluar funciones en un conjunto finito y manejable de puntos de
almacenamiento. Se define para ello una grilla de referencia de
\(N_s=9\) puntos equiespaciados por presa,
\begin{equation}\protect\phantomsection\label{eq-grilla-estado}{
\mathcal S_i^{\text{grid}} := \Big\{0,\ \tfrac{V_{i,\max}}{8},\ \tfrac{2V_{i,\max}}{8},\ \dots,\ V_{i,\max}\Big\}, \qquad i=1,2,
}\end{equation} con paso \(V_{1,\max}/8=2{,}025\) Mm³ y
\(V_{2,\max}/8=1{,}200\) Mm³ respectivamente. Esta grilla \textbf{no es
una discretización del MDP}: es un dispositivo de muestreo y de cómputo;
la política final se evalúa siempre en el valor continuo real del
almacenamiento.

\subsection{\texorpdfstring{Espacio de acciones
\((\mathcal{A})\)}{Espacio de acciones (\textbackslash mathcal\{A\})}}\label{sec-espacio-acciones}

Para cada etapa \(m\) se define una unidad de extracción de agua
\(u_m = (100,200,300,300,400,250)\) Mm³, y un conjunto de índices
factibles \(\mathcal K_{i,m}\) heredado de la grilla de decisión de la
PDE: \[
\begin{aligned}
\mathcal{K}_{0} &= \{0,\dots,20\}, &
\mathcal{K}_{1} &= \{0,\dots,15\}, &
\mathcal{K}_{2} &= \{0,\dots,7\}, \\
\mathcal{K}_{3} &= \{0,\dots,9\}, &
\mathcal{K}_{4} &= \mathcal{K}_{5} = \{0,\dots,12\}.
\end{aligned}
\] (idéntico para ambas presas en cada etapa). Estos índices \textbf{no
enumeran directamente el conjunto de acciones factibles del FQI}: se
usan únicamente para obtener la cota superior del turbinado factible en
la etapa, \begin{equation}\protect\phantomsection\label{eq-accion-max}{
A_{\max}(m) := u_m \cdot \max(\mathcal K_{m}).
}\end{equation} Por ejemplo, para \(m=0\): \(u_0=100\),
\(\max(\mathcal K_0)=20\), de modo que \(A_{\max}(0)=2{,}000\) Mm³ (así,
un par de índices como \((k_1,k_2)=(3,5)\) de la PDE correspondería a
turbinar \(300\) y \(500\) Mm³ respectivamente ---dentro de la cota---,
pero el FQI no restringe la acción a estos múltiplos enteros de
\(u_m\)).

El \textbf{espacio de acciones factible} para el MDP continuo es el
intervalo
\begin{equation}\protect\phantomsection\label{eq-espacio-acciones}{
\mathcal{A}_m := [0,A_{\max}(m)]\times[0,A_{\max}(m)] = \{a=(a_1,a_2)\}
}\end{equation} con \(a_i\) el volumen turbinado en la presa \(i\),
continuo. Al igual que con el estado, para fines de muestreo (Sec.
Section~\ref{sec-conjunto-aprendizaje}) y de maximización aproximada
(Sec. Section~\ref{sec-max-operador}) este intervalo se discretiza en
una grilla gruesa de \(N_a=9\) puntos equiespaciados,
\begin{equation}\protect\phantomsection\label{eq-grilla-accion}{
\mathcal A_m^{\text{grid}} := \Big\{0,\ \tfrac{A_{\max}(m)}{8},\ \dots,\ A_{\max}(m)\Big\}, \qquad |\mathcal A_m^{\text{grid}}|=9.
}\end{equation}

\begin{longtable}[]{@{}llll@{}}
\caption{Capacidad máxima de turbinado por etapa, calculada según
Equation~\ref{eq-accion-max}.}\label{tbl-accion-max}\tabularnewline
\toprule\noalign{}
Etapa \(m\) & \(u_m\) (Mm³) & \(\max(\mathcal K_m)\) & \(A_{\max}(m)\)
(Mm³) \\
\midrule\noalign{}
\endfirsthead
\toprule\noalign{}
Etapa \(m\) & \(u_m\) (Mm³) & \(\max(\mathcal K_m)\) & \(A_{\max}(m)\)
(Mm³) \\
\midrule\noalign{}
\endhead
\bottomrule\noalign{}
\endlastfoot
0 & 100 & 20 & 2,000 \\
1 & 200 & 15 & 3,000 \\
2 & 300 & 7 & 2,100 \\
3 & 300 & 9 & 2,700 \\
4 & 400 & 12 & 4,800 \\
5 & 250 & 12 & 3,000 \\
\end{longtable}

Los valores de \(u_m\) y \(\mathcal K_{m}\) se seleccionaron, en el
diseño original de la PDE, para que \(A_{\max}(m)\) no exceda la
capacidad instalada de las plantas ni los caudales históricos máximos
registrados; esa cota se conserva íntegramente en el modelo continuo de
FQI.

\subsection{\texorpdfstring{Dinámica de transición
\((\mathbb{P})\)}{Dinámica de transición (\textbackslash mathbb\{P\})}}\label{dinuxe1mica-de-transiciuxf3n-mathbbp}

\textbf{Balance de masa a nivel de etapa (usado para el conjunto de
aprendizaje).} Para construir las transiciones de entrenamiento (Sec.
Section~\ref{sec-conjunto-aprendizaje}), la afluencia se agrega a nivel
de etapa: sea
\begin{equation}\protect\phantomsection\label{eq-afluencia-etapa}{
\xi_{i,t,m} := \sum_{q\,:\,\mu(q)=m} W_{i,t,q}
}\end{equation} la afluencia histórica del año \(t\) agregada sobre las
quincenas de la etapa \(m\) (para la presa \(i\)), donde \(W_{i,t,q}\)
es el afluente neto observado en la quincena \(q\) del año \(t\) (dato
histórico). El balance de masa teórico (sin restricción de capacidad) al
aplicar la acción \(a_i\in\mathcal A_m^{\text{grid}}\) desde un estado
de la grilla \(s_i\in\mathcal S_i^{\text{grid}}\) es
\begin{equation}\protect\phantomsection\label{eq-transicion}{
\tilde s_i := s_i + \xi_{i,t,m} - a_i, \qquad i=1,2.
}\end{equation}

\textbf{Balance de masa a nivel de quincena (usado en el despliegue de
la política).} Al simular la política ya aprendida sobre la trayectoria
histórica completa (Sec. Section~\ref{sec-simulacion}), la decisión se
toma en cada quincena \(q\), usando el afluente histórico de esa
\textbf{misma} quincena:
\begin{equation}\protect\phantomsection\label{eq-transicion-quincenal}{
V_{i,t,q+1} = V_{i,t,q} + W_{i,t,q} - a_i, \qquad q=0,\dots,22,
}\end{equation} y, al cerrar el año (\(q=23\)), la transición conecta
con el primer registro del año siguiente,
\begin{equation}\protect\phantomsection\label{eq-transicion23}{
V_{i,t+1,0} = V_{i,t,23} + W_{i,t+1,0} - a_i.
}\end{equation} En ambos casos \(a_i\) es el volumen turbinado decidido
para la etapa \(m=\mu(q)\) vigente en esa quincena (Sec.
Section~\ref{sec-espacio-acciones}); \(a_i\) no cambia de valor dentro
de la misma quincena, pero sí puede cambiar de una quincena a otra
dentro de la misma etapa, porque la política se reconsulta con el estado
continuo corriente en cada quincena.

\textbf{Proyección al conjunto factible.} En ambos regímenes, el volumen
teórico \(\tilde s_i\) (o \(V_{i,t,q+1}\)) puede exceder \(V_{i,\max}\)
o resultar negativo. El volumen físicamente realizable se obtiene
proyectando sobre el intervalo factible,
\begin{equation}\protect\phantomsection\label{eq-proyeccion}{
s_i' := \operatorname{proj}_{[0,V_{i,\max}]}(\tilde s_i) = \min\big(\max(\tilde s_i,0),\,V_{i,\max}\big),
}\end{equation} donde \(\operatorname{proj}_{[0,V_{i,\max}]}\) denota al
operador que da el punto del intervalo \([0,V_{i,\max}]\) más cercano a
\(\tilde s_i\), y las magnitudes de exceso y faltante se registran como
derrame y déficit, respectivamente,
\begin{equation}\protect\phantomsection\label{eq-derrame-deficit}{
\text{derrame}_i := (\tilde s_i - V_{i,\max})^+, \qquad \text{déficit}_i := (-\tilde s_i)^+,
}\end{equation} con \(x^+:=\max(x,0)\).

La etapa evoluciona de manera determinista y cíclica,
\(m'=(m+1)\bmod M\), y el estado siguiente es \(s'=(s_1',s_2',m')\).

Desde la perspectiva de Puterman (1994), la dinámica del MDP se resume
en el núcleo de transición \(\mathbb P(\cdot\,|\,s,a)\), que en este
modelo \textbf{no se especifica en forma analítica cerrada}: toda la
incertidumbre proviene de la afluencia histórica \(\xi_{i,t,m}\) (o
\(W_{i,t,q}\)), tratada como una variable aleatoria cuya ley se aproxima
por la distribución empírica de las \(T=50\) realizaciones históricas
disponibles para cada etapa ---no por una tabla de probabilidades de
transición sobre estados discretos, puesto que \(\mathcal S\) es ahora
un continuo---.

\subsection{\texorpdfstring{Función de recompensa
\((\mathcal{R})\)}{Función de recompensa (\textbackslash mathcal\{R\})}}\label{sec-recompensa}

\subsubsection{Modelo hidráulico-energético
linealizado}\label{modelo-hidruxe1ulico-energuxe9tico-linealizado}

La energía generada depende del volumen turbinado y de una aproximación
lineal de la carga hidráulica. Para cada presa \(i\) se define la
función de carga \textbf{continua}
\begin{equation}\protect\phantomsection\label{eq-altura}{
\begin{aligned}
H_i(s_i,s_i')
&= H_{i,\max}\cdot
\frac{s_i+s_i'}{2\,V_{i,\max}}, \\[4pt]
H_{1,\max} &= 120\ \text{m (Angostura)}, \qquad
H_{2,\max} = 100\ \text{m (Malpaso)}.
\end{aligned}
}\end{equation} Es decir, la carga varía linealmente entre \(0\) m
(presa vacía) y la carga máxima propia de cada presa cuando ésta alcanza
su capacidad plena, y se aproxima la carga efectiva durante la
etapa/quincena mediante el promedio de la carga al inicio (\(s_i\)) y al
final (\(s_i'\)) del intervalo.

Sea \(\eta = 0.9\) la eficiencia de conversión hidráulica-eléctrica y
\(g = 9.81\) m/s² la aceleración de la gravedad. La energía generada por
la presa \(i\) al turbinar \(a_i\) (Mm³) bajo la carga \(H_i(s_i,s_i')\)
es \begin{equation}\protect\phantomsection\label{eq-energia}{
E_i(s,a,s') = \frac{\eta\, g\, H_i(s_i,s_i')\,(a_i\times 10^6)}{3600\times 10^6} \quad [\text{GWh}],
}\end{equation} donde \(10^6\) convierte Mm³ a m³ y \(3600\times10^6\)
convierte julios-equivalentes a GWh (con la densidad del agua
\(\rho=1000\) kg/m³ absorbida algebraicamente en la constante \(3600\)).
La energía total generada por el sistema en la etapa/quincena es
\begin{equation}\protect\phantomsection\label{eq-energia-total}{
E = E_1+E_2.
}\end{equation}

\subsubsection{Estructura de
penalizaciones}\label{estructura-de-penalizaciones}

La recompensa incorpora penalizaciones lineales, con coeficientes
calibrados para actuar como restricciones blandas de tipo \emph{big-M}
(su magnitud domina ampliamente el valor energético del agua, ver
ejemplo numérico más abajo). Se definen, por presa \(i\), las siguientes
tres penalizaciones:

\begin{itemize}
\item
  \textbf{Derrame} (excedente sobre la capacidad máxima):
  \begin{equation}\protect\phantomsection\label{eq-penalizacion-derrame}{
  \Pi_i^{\text{derr}} = c_{\text{derr}}\cdot(\tilde s_i - V_{i,\max})^+, \qquad c_{\text{derr}} := \frac{C_{\text{derr}}}{\Delta}.
  }\end{equation}
\item
  \textbf{Déficit} (faltante respecto de un balance no negativo).
  \emph{Proporcional a la magnitud del déficit, no una constante fija.}
  \begin{equation}\protect\phantomsection\label{eq-penalizacion-deficit}{
  \Pi_i^{\text{def}} = c_{\text{def}}\cdot(-\tilde s_i)^+, \qquad c_{\text{def}} := \frac{C_{\text{def}}}{\Delta}.
  }\end{equation}
\item
  \textbf{Violación de la curva guía}. \emph{Se penaliza cuando el
  almacenamiento proyectado \textbf{excede} la curva guía (un techo
  operativo de resguardo).}
  \begin{equation}\protect\phantomsection\label{eq-penalizacion-curva-guia}{
  \Pi_i^{\text{CG}} = c_{\text{CG}}\cdot\big(s_i' - \text{CG}_i(m)\big)^+, \qquad c_{\text{CG}} := \frac{C_{\text{CG}}}{\Delta}.
  }\end{equation}
\end{itemize}

con \(C_{\text{derr}}=C_{\text{def}}=10{,}000\),
\(C_{\text{CG}}=5{,}000\), de modo que
\(c_{\text{derr}}=c_{\text{def}}\approx16.67\) y
\(c_{\text{CG}}\approx8.33\) por Mm³. A modo de referencia de escala:
turbinar 1 Mm³ bajo una carga representativa de 100 m genera, por
Equation~\ref{eq-energia}, apenas \(\approx0.245\) GWh ---es decir,
\(c_{\text{derr}}\) es unas \(68\) veces mayor que el valor energético
marginal del mismo volumen, confirmando su papel de penalización
dominante y no de costo económico real.

\(\text{CG}_i(m)\) es la curva guía de la presa \(i\) en la etapa \(m\),
definida originalmente en unidades de índice,
\(\text{CG}_1^{\text{idx}}=(25,22,24,25,26,27)\),
\(\text{CG}_2^{\text{idx}}=(16,15,14,14,15,16)\), y convertida a volumen
mediante
\begin{equation}\protect\phantomsection\label{eq-curva-guia-vol}{
\text{CG}_i(m) := \Delta\cdot \text{CG}_i^{\text{idx}}(m).
}\end{equation}

La función de recompensa total, sumada sobre ambas presas, es
\begin{equation}\protect\phantomsection\label{eq-recompensa}{
\mathcal{R}(s,a,s') = \sum_{i=1}^{2}\Big[E_i - \Pi_i^{\text{derr}} - \Pi_i^{\text{def}} - \Pi_i^{\text{CG}}\Big].
}\end{equation}

En el presente trabajo no se modela una distribución analítica explícita
de la afluencia; en su lugar se construye un conjunto de aprendizaje
empírico, descrito con precisión en la Sec.
Section~\ref{sec-conjunto-aprendizaje} ---y que, como ahí se detalla,
\textbf{no} consiste en recorrer trayectorias históricas de decisiones
reales, sino en evaluar la dinámica y la recompensa sobre una malla
exhaustiva de pares estado-acción contra la totalidad de las
realizaciones históricas de afluencia. Este enfoque se alinea con un
esquema de \textbf{\emph{Batch Reinforcement Learning}} ``parcialmente
libre de modelo'' (Castelletti et al. (2010); Ernst et al. (2005)).

\section{Ecuación de optimalidad de
Bellman}\label{ecuaciuxf3n-de-optimalidad-de-bellman}

Se adopta la notación \(Q^*\) para la función de valor óptima, por ser
función del par estado-acción \((s,a)\) ---lo que la literatura de
aprendizaje por refuerzo llama función de valor-acción--- y por
consistencia con el nombre de la variable en el código (\texttt{Q}).

Conforme a la teoría clásica de MDPs descontados (Puterman (1994),
Teorema 6.2.3), existe una única función de valor-acción óptima \(Q^*\)
que satisface la ecuación de optimalidad de Bellman. La razón intuitiva
es el factor de descuento \(\gamma=0.95<1\): la influencia de una
decisión de hoy sobre lo que ocurre dentro de, digamos, 50 etapas está
multiplicada por \(\gamma^{50}\approx0.08\), y sigue decayendo
geométricamente cuanto más lejos se mire hacia el futuro. Ese
decaimiento garantiza que, sin importar con qué aproximación inicial se
empiece ---incluso \(Q^{(0)}\equiv0\), como hace el algoritmo (Sec.
Section~\ref{sec-conjunto-aprendizaje} y siguientes)---, al aplicar la
ecuación de Bellman una y otra vez las sucesivas aproximaciones se
acercan cada vez más entre sí, hasta converger a una única \(Q^*\). Esto
es válido tanto si \(\mathcal S\) es finito (el caso clásico de
Puterman) como si tiene una componente continua, como en este modelo
(Sec. Section~\ref{sec-espacio-estados}). \emph{(Formalmente, esta
propiedad se conoce como que el operador de Bellman es una contracción,
y su justificación rigurosa se apoya en el teorema del punto fijo de
Banach; no se desarrolla aquí por no ser el objeto central de este
capítulo.)}

Dado que la incertidumbre proviene de la afluencia histórica
\(\xi=(\xi_1,\xi_2)\), cuya ley se aproxima por su distribución empírica
\(\widehat P_m\) sobre las \(T=50\) realizaciones históricas de la etapa
\(m\) (Sec. Section~\ref{sec-conjunto-aprendizaje}), la ecuación de
optimalidad se escribe como una esperanza sobre dicha distribución
empírica ---y no como una suma sobre un conjunto discreto de estados
siguientes, formalismo que ya no aplica al ser \(\mathcal S\) un
continuo---: \begin{equation}\protect\phantomsection\label{eq-bellman}{
Q^{*}(s,a) = \mathbb{E}_{\xi\sim\widehat P_m}\Big[\,\mathcal R(s,a,s') + \gamma\max_{a'\in\mathcal A_{m'}} Q^{*}(s',a')\,\Big], \quad s' = f(s,a,\xi),
}\end{equation} con \(f\) la función de transición y proyección de las
ecuaciones Equation~\ref{eq-transicion}--Equation~\ref{eq-proyeccion}, y
\(m'=(m+1)\bmod M\).

Dado que no se dispone de \(\widehat P_m\) en forma cerrada sino de sus
\(T\) realizaciones muestrales, el algoritmo de FQI (Sec. siguiente)
trabaja con la versión \textbf{por muestra} de esta ecuación ---el
objetivo empírico de Bellman, Equation~\ref{eq-objetivo-bellman}---,
evaluada sobre cada una de las transiciones del conjunto de aprendizaje,
y no con la esperanza poblacional (Equation~\ref{eq-bellman})
directamente.

\section{Aproximación numérica mediante Fitted Q-Iteration
(FQI)}\label{aproximaciuxf3n-numuxe9rica-mediante-fitted-q-iteration-fqi}

La Programación Dinámica Estocástica está afectada por la maldición de
la dimensionalidad: el costo computacional crece exponencialmente con el
número de variables de estado, y requiere además un modelo de transición
explícito (Castelletti et al. (2010)). Por ello, el presente trabajo
aproxima \(Q^*\) mediante Fitted Q-Iteration, un algoritmo de
aprendizaje por refuerzo \emph{batch} que ajusta iterativamente un
aproximador funcional supervisado a partir de un conjunto de
transiciones muestreadas.

\subsection{Conjunto de aprendizaje ``parcialmente libre de
modelo''}\label{sec-conjunto-aprendizaje}

Concretamente, el conjunto de aprendizaje
\(\mathcal F=\{(x_j, r_j, s'_{1,j}, s'_{2,j}, m'_j)\}_{j=1}^N\) se
construye enumerando el producto cartesiano de: la grilla de estado
Equation~\ref{eq-grilla-estado}, la grilla de acción
Equation~\ref{eq-grilla-accion}, las \(M=6\) etapas, y las \(T=50\)
realizaciones históricas de afluencia por etapa
Equation~\ref{eq-afluencia-etapa}. El vector de características de cada
muestra es
\begin{equation}\protect\phantomsection\label{eq-caracteristicas}{
x_j = (s_{1,j}, s_{2,j}, m_j, a_{1,j}, a_{2,j}) \in \mathcal S_1^{\text{grid}}\times\mathcal S_2^{\text{grid}}\times\{0,\dots,5\}\times\mathcal A^{\text{grid}}_{m_j}\times\mathcal A^{\text{grid}}_{m_j},
}\end{equation} y, para cada combinación de estado-acción-etapa, se
generan \(T=50\) muestras ---una por cada año histórico \(t\)---
aplicando Equation~\ref{eq-transicion}, Equation~\ref{eq-proyeccion} y
Equation~\ref{eq-recompensa} con la afluencia real \(\Xi_{i,t,m_j}\) de
ese año. El tamaño total del conjunto de aprendizaje es
\begin{equation}\protect\phantomsection\label{eq-tamano-N}{
N = M\times|\mathcal S_1^{\text{grid}}|\times|\mathcal S_2^{\text{grid}}|\times|\mathcal A_1^{\text{grid}}|\times|\mathcal A_2^{\text{grid}}|\times T = 6\times9\times9\times9\times9\times50 = 1{,}968{,}300.
}\end{equation}

Esta receta de muestreo ---barrer exhaustivamente estado y acción sobre
una grilla, pero conservar \emph{todas} las realizaciones históricas de
la perturbación en cada punto, sin colapsarlas a una tabla de
probabilidades ni ajustarles una distribución paramétrica--- es la que
Castelletti et al. (2010) denominan \textbf{``parcialmente libre de
modelo''} (\emph{partial model-free}): la dinámica del balance de masa
sí se modela analíticamente, pero la ley de la afluencia se deja no
paramétrica.

\subsection{Iteración de Bellman
empírica}\label{iteraciuxf3n-de-bellman-empuxedrica}

Se inicializa \(Q^{(0)}(s,a) \equiv 0\) para todo \((s,a)\). Para cada
transición \((x_j,r_j,s'_{1,j},s'_{2,j},m'_j)\in\mathcal F\) del
conjunto de aprendizaje, se define el \textbf{objetivo de Bellman} de la
iteración \(n\) por
\begin{equation}\protect\phantomsection\label{eq-objetivo-bellman}{
y^{(n)}_j := r_j + \gamma \max_{a'\in\mathcal A_{m_j'}} Q^{(n-1)}(s'_{1,j},s'_{2,j},m_j',a'),
}\end{equation} que estima el valor-acción óptimo de \(x_j\) combinando
la recompensa observada \(r_j\) y el mejor valor alcanzable desde el
estado siguiente, según la aproximación de la iteración anterior. El
cómputo del máximo sobre \(a'\), dado que \((s'_{1,j},s'_{2,j})\) es en
general un punto continuo que no coincide con la grilla
Equation~\ref{eq-grilla-estado}, se describe en la Sec.
Section~\ref{sec-max-operador}.

\subsection{Aproximación del operador de maximización sobre estados
continuos}\label{sec-max-operador}

**Esta subsección resuelve un punto que el texto anterior dejaba
implícito sin justificación: cómo evaluar \(\max_{a'}Q^{(n-1)}(s',a')\)
cuando \(s'\) es continuo. Evaluarlo por fuerza bruta para cada una de
las \(N\approx 2\times10^6\) muestras sería prohibitivo, así que el
algoritmo precomputa, para cada etapa \(m\), una superficie de valor
máximo sobre la grilla de estado y la interpola:

\begin{enumerate}
\def\labelenumi{\arabic{enumi}.}
\item
  \textbf{Evaluación sobre grilla.} Para cada etapa \(m\), se evalúa
  \(Q^{(n-1)}\) en todo el producto
  \(\mathcal S_1^{\text{grid}}\times\mathcal S_2^{\text{grid}}\times\{m\}\times\mathcal A_1^{\text{grid}}(m)\times\mathcal A_2^{\text{grid}}(m)\),
  y se maximiza sobre las dos coordenadas de acción:
  \begin{equation}\protect\phantomsection\label{eq-maxq-grilla}{
  \widehat V^{(n-1)}(s_1^{(p)},s_2^{(q)},m) := \max_{a_1,a_2\in\mathcal A_m^{\text{grid}}} Q^{(n-1)}(s_1^{(p)},s_2^{(q)},m,a_1,a_2), \qquad p,q=1,\dots,9,
  }\end{equation} obteniéndose una matriz \(9\times9\) por etapa.
\item
  \textbf{Interpolación bilineal.} Sobre esta matriz se construye un
  interpolador continuo
  \(\widetilde V^{(n-1)}(\cdot,\cdot\,;m):\mathbb R^2\to\mathbb R\)
  (interpolación multilineal sobre grilla regular), con extrapolación
  lineal fuera de \([0,V_{1,\max}]\times[0,V_{2,\max}]\).
\item
  \textbf{Sustitución en el objetivo de Bellman.} El máximo de
  Equation~\ref{eq-objetivo-bellman} se aproxima evaluando el
  interpolador de la etapa siguiente en el punto continuo genuinamente
  alcanzado:
  \begin{equation}\protect\phantomsection\label{eq-maxq-interp}{
  \max_{a'\in\mathcal A_{m_j'}} Q^{(n-1)}(s'_{1,j},s'_{2,j},m_j',a') \ \approx\ \widetilde V^{(n-1)}\big(s'_{1,j},s'_{2,j};\,m_j'\big).
  }\end{equation}
\end{enumerate}

Esta aproximación reduce el costo de
\(O(N\times|\mathcal A^{\text{grid}}|^2)\) evaluaciones del regresor a
\(O\big(M\times|\mathcal S^{\text{grid}}|^2\times|\mathcal A^{\text{grid}}|^2 + N\big)\),
al reutilizar la misma superficie interpolada para las \(N/M\) muestras
de cada etapa.

\subsection{Aproximación funcional por regresión no paramétrica
(ExtraTrees)}\label{aproximaciuxf3n-funcional-por-regresiuxf3n-no-paramuxe9trica-extratrees}

Una vez calculados los objetivos de Bellman \(y^{(n)}\)
(Equation~\ref{eq-objetivo-bellman}, con el máximo aproximado según
Equation~\ref{eq-maxq-interp}), se ajusta un aproximador \(Q^{(n)}\) que
minimiza el error cuadrático entre las predicciones \(Q^{(n)}(x_j)\) y
los objetivos \(y^{(n)}_j\) sobre todas las transiciones del conjunto de
aprendizaje. Se utiliza, como propone Ernst et al. (2005), un ensamble
de \emph{ExtraTrees} (árboles extremadamente aleatorizados), adecuado
para aproximar funciones de valor en espacios continuos.

\textbf{Hiperparámetros (corregidos respecto de la versión previa):}

\begin{longtable}[]{@{}llll@{}}
\caption{Hiperparámetros del ensamble de Extra-Trees, según el código,
frente a lo reportado en la versión previa del
capítulo.}\label{tbl-hiperparametros}\tabularnewline
\toprule\noalign{}
Hiperparámetro & Símbolo & Valor implementado & \\
\midrule\noalign{}
\endfirsthead
\toprule\noalign{}
Hiperparámetro & Símbolo & Valor implementado & \\
\midrule\noalign{}
\endhead
\bottomrule\noalign{}
\endlastfoot
Número de árboles & \(M_{\text{tree}}\) & 300 & \\
Profundidad máxima & --- & \textbf{sin límite}
(\texttt{max\_depth=None}) & \\
Muestras mínimas por hoja & \(n_{\min}\) & \textbf{50} (\(=T\)) & \\
Variables candidatas por corte & \(K\) & \(n=5\)
(\texttt{max\_features=1.0}) & \\
\end{longtable}

El valor \(n_{\min}=T=50\) no es arbitrario: al exigir que cada hoja del
árbol contenga al menos tantas muestras como años históricos de
afluencia, se favorece que cada hoja agrupe una fracción representativa
de las \(T\) realizaciones de \(\xi\) asociadas a un mismo \((s,a)\)
aproximado, de modo que el promedio dentro de la hoja aproxime, por la
ley (empírica) de los grandes números, la esperanza condicional
\(\mathbb E_{\xi\sim\widehat P_m}[\cdot]\) que aparece en
Equation~\ref{eq-bellman}. Cada árbol particiona el espacio
estado-acción en regiones (hojas) mediante cortes con umbral aleatorio
(no exhaustivo, a diferencia de CART/Random Forest) sobre las \(K=n=5\)
variables de entrada, y predice, para un punto de consulta, el promedio
del objetivo de regresión entre las muestras de entrenamiento de su
misma hoja. La predicción del ensamble es el promedio de las
\(M_{\text{tree}}=300\) predicciones individuales.

\begin{quote}
\textbf{Nota sobre la Figura de árbol de decisión.} La versión previa
incluía una figura con la leyenda \emph{``Árbol de decisión con
profundidad máxima de 14 y un mínimo de 5 muestras por hoja''}. Dado que
el código real no limita la profundidad (\texttt{max\_depth=None}) y usa
\(n_{\min}=50\), esa leyenda ---y, en la práctica, la figura misma, que
mostraría un árbol truncado no representativo--- debe corregirse o
sustituirse por una ilustración esquemática que no afirme una
profundidad máxima inexistente.
\end{quote}

El algoritmo repite este procedimiento iterativamente.

\subsection{Criterio de paro basado en desempeño
simulado}\label{sec-criterio-paro}

Un criterio de convergencia del tipo punto-fijo ---detener cuando
\(\|Q^{(n)}-Q^{(n-1)}\|\) cae bajo una tolerancia--- \textbf{no es
aplicable} cuando el aproximador es un ensamble de Extra-Trees: la
aleatorización propia de la construcción de cada árbol (elección
aleatoria de umbrales de corte) hace que \(Q^{(n)}\) y \(Q^{(n-1)}\)
difieran por ruido de ajuste incluso si ya se hubiese alcanzado, en
esencia, la función de valor-acción óptima; esta distancia no converge a
cero. Así lo argumentan explícitamente Castelletti et al. (2010), y así
lo implementa el código, que en su lugar monitorea el \textbf{desempeño
operativo simulado} de la política inducida en cada iteración.

\textbf{Simulación de la política golosa.} Cada \(\Delta h\) iteraciones
(en el código, \(\Delta h=4\)), se simula la operación del sistema bajo
la política golosa de \(Q^{(n)}\) sobre la trayectoria histórica
completa, en resolución de quincena (Sec. Section~\ref{sec-simulacion}),
obteniendo energías \(e^{(i)}_{t,q}\), derrames \(\delta^{(i)}_{t,q}\) y
déficits \(\epsilon^{(i)}_{t,q}\) por presa.

\textbf{Funcional de desempeño.} Se define
\begin{equation}\protect\phantomsection\label{eq-desempeno-Jn}{
\widehat J_n := \sum_{t,q}\big(e^{(1)}_{t,q}+e^{(2)}_{t,q}\big) - c_{\text{derr}}\sum_{t,q}\big(\delta^{(1)}_{t,q}+\delta^{(2)}_{t,q}\big) - c_{\text{def}}\sum_{t,q}\big(\epsilon^{(1)}_{t,q}+\epsilon^{(2)}_{t,q}\big),
}\end{equation} una medida de desempeño agregada, \textbf{no descontada}
y evaluada con afluencias reales quincenales ---distinta, por tanto, del
objetivo de ajuste Equation~\ref{eq-objetivo-bellman}, que opera a nivel
de una transición individual y con descuento \(\gamma\)---.

\textbf{Regla de paro.} Se detiene la iteración la primera vez que se
observa una caída de \(\widehat J_n\) respecto del checkpoint
inmediatamente anterior (una ``reversión''):
\begin{equation}\protect\phantomsection\label{eq-criterio-paro}{
n^\star := \min\{\, n-\Delta h \ :\ \widehat J_{n} < \widehat J_{n-\Delta h} \,\},
}\end{equation} interpretando la reversión como evidencia de que, a
partir de ese punto, predominan las fluctuaciones aleatorias del ajuste
sobre el aprendizaje genuino (Castelletti et al. (2010)). Si no se
observa ninguna reversión dentro del número máximo de iteraciones
permitido, se usa como respaldo el mejor checkpoint observado,
\(n^\star:=\arg\max_n \widehat J_n\). La aproximación final es
\(Q^{*}:=Q^{(n^\star)}\), recuperada de una caché de modelos
serializados guardada en cada punto de verificación (necesaria porque,
por la aleatorización del ensamble, \(Q^{(n^\star)}\) no podría
reproducirse exactamente reentrenando desde cero).

\section{Política óptima y simulación operativa}\label{sec-simulacion}

Una vez obtenida \(Q^{*}\), se deriva la política óptima (golosa)
\(\pi^{*}:\mathcal S\to\mathcal A\),
\begin{equation}\protect\phantomsection\label{eq-politica-optima}{
\pi^{*}(s) := \operatorname*{arg\,max}_{a\in\mathcal A_m^{\text{grid}}} Q^{*}(s,a),
}\end{equation} donde la maximización se realiza, en la práctica,
mediante búsqueda exhaustiva sobre la grilla de acción de la etapa
vigente Equation~\ref{eq-grilla-accion} (no sobre el conjunto continuo
\(\mathcal A_m\)), evaluando \(Q^{*}\) en el estado continuo real ---no
discretizado--- del sistema.

La política se valida mediante simulación en lazo cerrado sobre la serie
histórica completa, siguiendo la práctica estándar de validación de
políticas de control aprendidas documentada en Bertsekas (2019). El
algoritmo, corregido para ser fiel al código, opera así:

\begin{enumerate}
\def\labelenumi{\arabic{enumi}.}
\item
  \textbf{Inicialización:} se fijan los volúmenes iniciales
  \(V_{1,0}=10{,}000\) Mm³, \(V_{2,0}=6{,}000\) Mm³ (valores del código;
  no se especificaban antes).
\item
  \textbf{Mapeo estacional:} para cada quincena \(q\) (recorriendo
  \(t=1,\dots,T\) y, dentro de cada año, \(q=0,\dots,23\)), se determina
  \(m=\mu(q)\) y la grilla de acción \(\mathcal A_m^{\text{grid}}\)
  vigente.
\item
  \st{\textbf{Discretización:} $s_i=\mathcal D(V_i,N_i)$} ---
  \textbf{paso eliminado}: el estado usado para consultar \(Q^*\) es el
  volumen continuo real \((V_{1,t,q}, V_{2,t,q})\), sin redondeo a
  índice alguno.
\item
  \textbf{Selección de acción}, por búsqueda exhaustiva sobre la grilla
  de acción de 9×9 combinaciones de la etapa vigente: \[
  (a_1^*,a_2^*) = \operatorname*{arg\,max}_{(a_1,a_2)\,\in\,\mathcal A_1^{\text{grid}}(m)\times\mathcal A_2^{\text{grid}}(m)} Q^{*}(V_{1,t,q},V_{2,t,q},m,a_1,a_2).
  \]
\item
  \textbf{Balance hídrico}, con el afluente de la \textbf{misma}
  quincena \(q\) (corrección respecto de la versión previa, que usaba
  \(q+1\)): \[
  V_{i,t,q+1} = V_{i,t,q} + W_{i,t,q} - a_i^*, \qquad i=1,2.
  \]
\item
  \textbf{Proyección física}, incluyendo tanto el derrame como el
  déficit (la versión previa omitía el déficit): \[
  V_i^{\text{new}} = \operatorname{proj}_{[0,V_{i,\max}]}(V_{i,t,q+1}), \qquad
  \text{derrame}_i = (V_{i,t,q+1}-V_{i,\max})^+, \qquad
  \text{déficit}_i = (-V_{i,t,q+1})^+.
  \]
\item
  \textbf{Cálculo energético}, con la carga continua y las cargas
  máximas distintas por presa (corrección de Equation~\ref{eq-altura}):
  \[
  H_i = H_{i,\max}\cdot\frac{V_{i,t,q}+V_i^{\text{new}}}{2\,V_{i,\max}}, \qquad
  E_i = \frac{\eta g\, H_i\,(a_i^*\times10^6)}{3600\times10^6}.
  \]
\item
  \textbf{Registro y actualización:} se almacenan energía, derrame,
  déficit y almacenamiento; el estado se actualiza a
  \(V_i^{\text{new}}\) para la siguiente quincena.
\end{enumerate}

\subsection{Métricas de desempeño y validación
estadística}\label{sec-metricas}

\begin{tcolorbox}[enhanced jigsaw, arc=.35mm, bottomrule=.15mm, bottomtitle=1mm, breakable, colback=white, colbacktitle=quarto-callout-warning-color!10!white, colframe=quarto-callout-warning-color-frame, coltitle=black, left=2mm, leftrule=.75mm, opacityback=0, opacitybacktitle=0.6, rightrule=.15mm, title=\textcolor{quarto-callout-warning-color}{\faExclamationTriangle}\hspace{0.5em}{Cifras pendientes de regenerar}, titlerule=0mm, toprule=.15mm, toptitle=1mm]

Las tablas siguientes (Table~\ref{tbl-desempeno-energetico},
Table~\ref{tbl-desempeno-hidrico}) y la conclusión que las acompañaba en
la versión previa reportan un derrame total del sistema de
\(\approx 34{,}053\) Mm³ y una lectura crítica sobre la mala gestión de
la política. Esta cifra \textbf{es incompatible con la propia validación
documentada en el encabezado del script corregido}, que reporta ``0 Mm³
de derrame y 0 Mm³ de déficit en ambas presas'' al ejecutar exactamente
el algoritmo aquí descrito. Es decir, las cifras de la versión previa
corresponden, con alta probabilidad, a una ejecución de una versión
anterior y no fiel del modelo (la que originó la auditoría documentada
en el encabezado de \texttt{FQI\_corregido.py}), no a la versión actual.
No se sustituyen aquí por cifras nuevas para no fabricar resultados: se
recomienda \textbf{volver a ejecutar el script corregido sobre los datos
reales} y regenerar ambas tablas, así como reescribir el párrafo de
interpretación en función de los resultados que se obtengan.

\end{tcolorbox}

\begin{longtable}[]{@{}
  >{\raggedright\arraybackslash}p{(\linewidth - 6\tabcolsep) * \real{0.2500}}
  >{\raggedright\arraybackslash}p{(\linewidth - 6\tabcolsep) * \real{0.2500}}
  >{\raggedright\arraybackslash}p{(\linewidth - 6\tabcolsep) * \real{0.2500}}
  >{\raggedright\arraybackslash}p{(\linewidth - 6\tabcolsep) * \real{0.2500}}@{}}
\caption{Resumen del desempeño energético del sistema bajo la política
óptima aprendida mediante FQI (plantilla; valores por regenerar con
\protect\texttt{simulacion\_FQI\_corregida.csv}).}\label{tbl-desempeno-energetico}\tabularnewline
\toprule\noalign{}
\begin{minipage}[b]{\linewidth}\raggedright
Métrica
\end{minipage} & \begin{minipage}[b]{\linewidth}\raggedright
Angostura
\end{minipage} & \begin{minipage}[b]{\linewidth}\raggedright
Malpaso
\end{minipage} & \begin{minipage}[b]{\linewidth}\raggedright
Sistema Total
\end{minipage} \\
\midrule\noalign{}
\endfirsthead
\toprule\noalign{}
\begin{minipage}[b]{\linewidth}\raggedright
Métrica
\end{minipage} & \begin{minipage}[b]{\linewidth}\raggedright
Angostura
\end{minipage} & \begin{minipage}[b]{\linewidth}\raggedright
Malpaso
\end{minipage} & \begin{minipage}[b]{\linewidth}\raggedright
Sistema Total
\end{minipage} \\
\midrule\noalign{}
\endhead
\bottomrule\noalign{}
\endlastfoot
Energía promedio quincenal {[}GWh{]} & \emph{(pendiente)} &
\emph{(pendiente)} & \emph{(pendiente)} \\
Energía anual estimada {[}GWh{]} & \emph{(pendiente)} &
\emph{(pendiente)} & \emph{(pendiente)} \\
Contribución al sistema {[}\%{]} & \emph{(pendiente)} &
\emph{(pendiente)} & 100.0 \\
\end{longtable}

\begin{longtable}[]{@{}llll@{}}
\caption{Resumen del desempeño hídrico del sistema bajo la política
óptima aprendida mediante FQI (plantilla; valores por
regenerar).}\label{tbl-desempeno-hidrico}\tabularnewline
\toprule\noalign{}
Métrica & Angostura & Malpaso & Sistema Total \\
\midrule\noalign{}
\endfirsthead
\toprule\noalign{}
Métrica & Angostura & Malpaso & Sistema Total \\
\midrule\noalign{}
\endhead
\bottomrule\noalign{}
\endlastfoot
Derrame total {[}Mm³{]} & \emph{(pendiente)} & \emph{(pendiente)} &
\emph{(pendiente)} \\
Déficit total {[}Mm³{]} & \emph{(pendiente)} & \emph{(pendiente)} &
\emph{(pendiente)} \\
\end{longtable}

Una vez regeneradas las tablas a partir de una corrida real del script,
la interpretación debe contrastar específicamente los totales de derrame
y déficit contra la validación reportada en el propio código (0 Mm³ en
ambos rubros) antes de emitir un juicio sobre la calidad de la política
aprendida; de mantenerse una discrepancia con lo documentado en el
script, ésta debería investigarse como posible diferencia entre el
entorno de ejecución (datos, semilla aleatoria, versión de librerías) y
no asumirse automáticamente como una falla de la política.

\section{Conclusiones del capítulo}\label{conclusiones-del-capuxedtulo}

\begin{itemize}
\tightlist
\item
  Síntesis de contribuciones matemáticas y numéricas.
\item
  Reafirmación de la coherencia entre formulación teórica,
  implementación algorítmica y validación operativa ---pendiente de
  completar una vez regeneradas las métricas de la Sec.
  Section~\ref{sec-metricas}---.
\item
  Declaración de cierre alineada con los objetivos de la tesis.
\end{itemize}

\bookmarksetup{startatroot}

\chapter{}\label{section}

\bookmarksetup{startatroot}

\chapter{}\label{section-1}

\bookmarksetup{startatroot}

\chapter*{Referencias}\label{referencias}
\addcontentsline{toc}{chapter}{Referencias}

\markboth{Referencias}{Referencias}

\protect\phantomsection\label{refs}
\begin{CSLReferences}{1}{1}
\bibitem[\citeproctext]{ref-Bertsekas2019}
Bertsekas, Dimitri P. 2019. \emph{Reinforcement Learning and Optimal
Control}. Athena Scientific.

\bibitem[\citeproctext]{ref-Castelletti2010}
Castelletti, A., S. Galelli, M. Restelli, and R. Soncini-Sessa. 2010.
{``Tree-Based Reinforcement Learning for Optimal Water Reservoir
Operation.''} \emph{Water Resources Research} 46: W09507.
\url{https://doi.org/10.1029/2009WR008898}.

\bibitem[\citeproctext]{ref-Ernst2005}
Ernst, Damien, Pierre Geurts, and Louis Wehenkel. 2005. {``Tree-Based
Batch Mode Reinforcement Learning.''} \emph{Journal of Machine Learning
Research} 6: 503--56.
\url{http://www.jmlr.org/papers/volume6/ernst05a/ernst05a.pdf}.

\bibitem[\citeproctext]{ref-ibmRL}
IBM. 2025. {``¿Qué es el aprendizaje por refuerzo?''} IBM.
\url{https://www.ibm.com/mx-es/think/topics/reinforcement-learning?mhsrc=ibmsearch_a&mhq=aprendizaje\%20por\%20refuerzo}.

\bibitem[\citeproctext]{ref-puterman1994markov}
Puterman, Martin L. 1994. \emph{Markov Decision Processes: Discrete
Stochastic Dynamic Programming}. Wiley Series in Probability and
Statistics. John Wiley \& Sons.
\url{https://doi.org/10.1002/9780470316887}.

\end{CSLReferences}




\end{document}
